\section{Discussion and Limitations}
\label{sec:discussion}

\subsection{Mechanism Interaction}

The three mechanisms are complementary, not redundant.  Input filtering
bounds fusion cost, which is the first stage to exceed the deadline as
$N$ grows.  Temporal amortization reduces prediction cost at any $N$ by
exploiting inter-tick coherence.  Risk budgeting allocates the
remaining prediction budget to the objects that matter most for planner
safety.  The ablation (\cref{tab:ablation_deadline}) shows that no
single mechanism achieves 100\% compliance across the full $N$ range;
the combination does.

\subsection{Generalization}

The cost model coefficients (\cref{eq:cost}) are hardware-specific
(calibrated on RTX~4080~SUPER).  Deploying on different MEC hardware
(e.g., NVIDIA Orin, T4) requires re-profiling per-stage latencies.
The controller structure and decision variables are hardware-agnostic:
the same $K$-search and risk-budgeted scheduling apply once the cost
model is re-calibrated.

The input contribution filter currently uses a BEV conflict-zone
integral.  Alternative scoring functions (e.g., information-theoretic
measures, learned attention) could replace it without changing the
controller's outer loop.

\subsection{Limitations}

\textbf{Offline evaluation only.}  All measurements are from an
offline profiler replaying recorded features.  Closed-loop CARLA
evaluation (where planner decisions feed back into vehicle behavior)
is in progress.

\textbf{Synthetic detection workload.}  Detections fed to the tracker
are synthetically generated from ground-truth trajectories, not from
the fusion model's actual output.  The timing measurements use real
features and real model inference, but the detection-to-tracking
interface is simulated.

\textbf{Divergence thresholds are heuristic.}  The velocity-normalized
threshold $v_i \cdot \tau$ (\cref{eq:divergence}) is manually tuned.
A learned threshold policy conditioned on object type and scene context
may improve amortization rates without sacrificing quality.

\textbf{Risk score not validated against closed-loop outcomes.}  The
multiplicative risk score (\cref{eq:risk}) is a surrogate calibrated
against offline prediction error, not against planner collision rates.
A learned risk function trained on closed-loop data would strengthen
the safety argument.

\textbf{Cache coherence.}  Tracker ID switches invalidate the
prediction cache.  The current implementation treats an ID switch as
a new object (cold start), which temporarily increases prediction
cost.  Integrating re-identification could mitigate the penalty.

\textbf{Dataset scale.}  OPV2V provides at most 2--5 real CAVs per
scene.  For $N > 5$, we replicate features to synthesize larger
agent counts.  Evaluation on Multi-V2X~\cite{multiv2x2024}
(up to 31 real CAVs per frame, 410 total CAVs, 56 RSUs) is in
progress and will provide genuine dense-scale workloads.

\subsection{Comparison with Prior Work}

CMP~\cite{cmp2024} is the closest system: it combines CoBEVT fusion,
AB3DMOT tracking, and MTR prediction.  CMP reports 67.3\,ms total
latency on an A6000 at 2--7 agents, with no deadline constraint and
no adaptation mechanism.  Our work extends the evaluation to $N{=}64$,
introduces the deadline cliff characterization, and proposes the
adaptive controller that maintains compliance at dense scale.

EMP~\cite{EMP} and Harbor~\cite{harbor} address edge-hosted
cooperative perception but stop at detection (no tracking or
prediction), report latency at small $N$, and do not model the
joint-pipeline deadline cliff.
